%%%
%%% Radical "⽘"
%%%
\section*{Radical 89: ``⽘''}\addcontentsline{toc}{section}{Radical 89: ⽘}\addcontentsline{loh}{figure}{\#\#\#\# 89: ⽘}

%%%%%%%%%% 爽 %%%%%%%%%%
\subsection*{爽}\addcontentsline{loh}{figure}{爽}

\begin{Entry}{爽}{11}{⽘}
  \begin{Phonetics}{爽}{shuang3}[][HSK 6]
    \definition{adj.}{claro; nítido; brilhante | franco; de coração aberto; direto | relaxado; confortável}
    \definition{v.}{desviar; afastar | tornar confortável; ficar confortável}
  \synonymref{畅}{chang4}
  \synonymref{快}{kuai4}
  \synonymref{朗}{lang3}
  \synonymref{凉}{liang2}
  \synonymref{明}{ming2}
  \synonymref{舒}{shu1}
  \synonymref{直}{zhi2}
  \antonymref{扭}{niu3}
  \antonymref{污}{wu1}
  \antonymref{脏}{zang1}
  \end{Phonetics}
\end{Entry}

\begin{Entry}{爽快}{11,7}{⽘,⼼}
  \begin{Phonetics}{爽快}{shuang3kuai5}[][HSK 7-9]
    \definition{adj.}{relaxado; confortável; devido a fatores ambientais e climáticos, sinto"-me relaxado e confortável tanto física quanto mentalmente | direto; franco; isso descreve alguém que consegue tomar decisões rapidamente, seja falando ou agindo, sem hesitação ou indecisão}
  \synonymref{干脆}{gan1cui4}
  \synonymref{坦率}{tan3shuai4}
  \antonymref{迟疑}{chi2yi2}
  \antonymref{犹豫}{you2yu4}
  \end{Phonetics}
\end{Entry}

%%%%% EOF %%%%%

